\%============================================================
\chapter{Bai Hoc: Loi Noi Va Hanh Dong --- Hanh Dong Noi To Hon Loi Noi (Actions Speak Louder Than Words)}
\begin{center}
  \textit{\color{truyengold}``Lam tot con hon la noi hay.''}\\[4pt]
  \textit{(Well done is better than well said.)}\\[4pt]
  \small\color{truyenblue}--- Co Kim Dong Tay ---
\end{center}
\ngancach

\section{Gandhi --- Hay La Su Thay Doi Ban Muon Thay}
\begin{truyen}{Gandhi --- Hay La Su Thay Doi Ban Muon Thay}{Mahatma Gandhi --- An Do TK 20}
\chuhoa{M}{}M ot ba me dan con trai den gap Gandhi va noi: 'Nha lanh dao oi, xin ong noi voi con toi dung an duong --- no an qua nhieu duong.'\\
\textit{(A mother brought her son to Gandhi and said: 'Leader, please tell my son not to eat sugar --- he eats too much.')}

Gandhi im lang roi noi: 'Tuan sau hay cho nguoi den gap toi.' Ba me boi roi nhung van lam theo.\\
\textit{(Gandhi was silent then said: 'Bring him back next week.' The mother was puzzled but complied.)}

Tuan sau, Gandhi chi nhin dua tre noi: 'Con oi, dung an qua nhieu duong --- no khong tot cho suc khoe.'\\
\textit{(The following week, Gandhi simply looked at the boy and said: 'Son, don't eat too much sugar --- it's not good for you.')}

Ba me hoi: 'Tai sao tuan truoc ong khong noi? Sao lai cho chung toi di roi den?' Gandhi tra loi:\\
\textit{(The mother asked: 'Why didn't you say that last week? Why make us go and come back?' Gandhi replied:)}

'Tuan truoc, chinh toi cung dang an nhieu duong. Toi can mot tuan de bo no truoc khi day nguoi khac.'\\
\textit{('Last week I myself was still eating much sugar. I needed a week to give it up before teaching others.')}

\end{truyen}
\begin{baihoc}
  Chi day dieu minh thuc hanh; khong ai co quyen day nguoi khac dieu ma minh chua song qua.\\
  \textit{(Only teach what you practice; no one has the right to teach others what they have not yet lived.)}
\end{baihoc}

\section{Lincoln Va Buc Thu Khong Gui}
\begin{truyen}{Lincoln Va Buc Thu Khong Gui}{Abraham Lincoln --- Hoa Ky 1864}
\chuhoa{K}{}K hi ai do lam Lincoln tuc gian, ong co mot thoi quen la viet ngay mot buc thu gian du vao ho.\\
\textit{(Whenever someone angered Lincoln, he had a habit of immediately writing a furious letter to them.)}

Ong viet that that, xuong het buc xuc, roi... gap thu lai de o ngan ban khong bao gio gui.\\
\textit{(He wrote everything out, releasing all frustration, then... folded the letter, kept it in his drawer, never mailed it.)}

Bi thu ky hoi tai sao, ong giai thich: 'Sau khi viet xong, toi da noi nhung dieu minh can noi.'\\
\textit{(Asked by his secretary why, he explained: 'After writing it, I've already said what I needed to say.')}

'Gui di thi chi gay them mau thuan --- ma toi can giai phap, khong can mot ke thu them.'\\
\textit{('Sending it would only create more conflict --- and I need solutions, not another enemy.')}

Sau khi mat ong, nguoi ta tim thay hang chuc buc thu gian du chua bao gio gui trong ngan ban cua ong.\\
\textit{(After his death, people found dozens of furious, never-sent letters in his desk drawers.)}

\end{truyen}
\begin{baihoc}
  Hanh dong khon ngoan nhat doi khi la khong hanh dong --- viet ra de xa stress nhung dung bao gio gui di khi dang tuc.\\
  \textit{(The wisest action is sometimes inaction --- write it out to release stress but never send it when angry.)}
\end{baihoc}

\section{Socrates Va Tam Cai Ro Cua Tin Tuc}
\begin{truyen}{Socrates Va Tam Cai Ro Cua Tin Tuc}{Socrates --- Hy Lap TK 5 TCN}
\chuhoa{M}{}M ot man nguoi chay den gap Socrates: 'Thay oi, toi can noi cho thay nghe mot chuyen ve ban cua thay!'\\
\textit{(A man ran to Socrates: 'Master, I need to tell you something about your friend!')}

Socrates noi: 'Khoan da --- truoc het ban hay cho loi noi cua minh qua ba cai ro.'\\
\textit{(Socrates said: 'Wait --- first let your words pass through three sieves.')}

'Cai ro thu nhat: No co that khong? Cai ro thu hai: Co tot khong? Cai ro thu ba: Co can thiet khong?'\\
\textit{('First sieve: Is it true? Second sieve: Is it good? Third sieve: Is it necessary?')}

Nguoi dan ong do lung: 'Cai ma toi dinh noi khong that chac, khong tot, va co le cung khong can thiet.'\\
\textit{(The man faltered: 'What I was going to say is not certainly true, not good, and perhaps not necessary either.')}

Socrates mim cuoi: 'Neu loi noi cua ban khong vượt qua duoc ba cai ro nay --- hay giu no lai cho minh.'\\
\textit{(Socrates smiled: 'If your words cannot pass these three sieves --- keep them to yourself.')}

\end{truyen}
\begin{baihoc}
  Truoc khi noi, hoi ban than: No co that? No co tot? No co can? Neu khong co it nhat hai cai, thi tot nhat la dung noi.\\
  \textit{(Before speaking, ask yourself: Is it true? Is it good? Is it needed? If you can't answer yes to at least two, it's best left unsaid.)}
\end{baihoc}

\section{Thomas Edison --- Lam Duc Ket Qua Bai Hoc}
\begin{truyen}{Thomas Edison --- Lam Duc Ket Qua Bai Hoc}{Thomas Edison --- Hoa Ky 1879}
\chuhoa{S}{}S au khi chay thi den dien 10.000 lan, ai do hoi Edison co cam thay minh that bai khong.\\
\textit{(After testing the light bulb 10,000 times, someone asked Edison if he felt like a failure.)}

Edison tra loi binh tinh: 'Toi chua that bai mot lan nao ca --- toi da hoc duoc 10.000 cach khong hoat dong.'\\
\textit{(Edison replied calmly: 'I have not failed even once --- I have learned 10,000 ways that do not work.')}

Nhung dieu Edison noi khong chi la cau noi dep; oa la cach on thuc su song va lam viec.\\
\textit{(What Edison said was not just a pretty quote; that was how he genuinely lived and worked.)}

Ong giu nhat ky thi nghiem --- ghi lai ket qua tung lan that bai mot cach chi tiet va khoa hoc.\\
\textit{(He kept experiment journals --- recording the results of each failed attempt in meticulous scientific detail.)}

Khi den dien cuoi cung sang len, no sang len tren nen tang cua chinh nhung lan that bai do.\\
\textit{(When the light bulb finally lit up, it did so built upon the foundation of those very failures.)}

\end{truyen}
\begin{baihoc}
  Nguoi goi bat cu ket qua nao la 'that bai' thuc ra dang tu choi bai hoc quy gia nhat trong thuc nghiem.\\
  \textit{(Anyone who calls any result 'failure' is actually refusing the most valuable lesson in the experiment.)}
\end{baihoc}

\section{Nguyen Trai --- Tu La Trung Va Nghia Quan}
\begin{truyen}{Nguyen Trai --- Tu La Trung Va Nghia Quan}{Nguyen Trai --- Viet Nam TK 15}
\chuhoa{K}{}K hi Nguyen Trai viet Binh Ngo Dai Cao, ong khong chi viet mot tuyen ngon chien thang --- ong viet no de day dan.\\
\textit{(When Nguyen Trai wrote the Great Proclamation of Victory over the Wu, he was not just writing a victory declaration --- he was teaching the people.)}

Ong viet: 'Viec nhan nghia cot o yen dan; quan dieu co du thuat la trong bi ham ho.'\\
\textit{(He wrote: 'The purpose of humanity and justice is to bring peace to the people; military strategy must prioritize removing disaster.')}

Bai hoc ong neu ra o day la: dung suc manh quan su nen phai di kem theo muc dich nhan van.\\
\textit{(The lesson he stated was: the use of military force must be accompanied by a humane purpose.)}

Sau khi duoi duoc giac Minh, ong de nghi tha bo tuong bi bat thay vi xu tham --- vi 'giang nhuc ke thua chua bao gio tao ra hoa binh dai lau'.\\
\textit{(After expelling the Ming invaders, he proposed releasing captured enemy generals instead of executing them --- because 'humiliating the defeated never creates lasting peace.')}

Chinh sach nay giup Viet Nam giu hoa binh vung chac voi trung Hoa trong su the ky tiep theo.\\
\textit{(This policy helped Vietnam maintain stable peace with China for the following century.)}

\end{truyen}
\begin{baihoc}
  Hanh dong nhan nghia sau khi chien thang con kho hon hanh dong anh hung trong khi chien tranh --- va cung quan trong hon nhieu.\\
  \textit{(Acting with humanity after victory is harder than acting heroically during war --- and far more important.)}
\end{baihoc}

\section{Marcus Aurelius --- De Quoc Rung Manh Nhung Tam Long Binh Than}
\begin{truyen}{Marcus Aurelius --- De Quoc Rung Manh Nhung Tam Long Binh Than}{Marcus Aurelius --- Rome TK 2}
\chuhoa{N}{}N guoi ta noi rang Marcus Aurelius la vi Hoang De Roman vi dai nhat vi ong la nha chien thuat tai ba.\\
\textit{(People say Marcus Aurelius was Rome's greatest emperor because he was a brilliant strategist.)}

Nhung bi mat thuc su nam o nhat ky ca nhan cua ong --- Suy Nghi -- ma ong khong bao gio dinh se xuat ban.\\
\textit{(But the real secret lay in his personal journal --- Meditations --- which he never intended to publish.)}

Ong viet cho chinh minh: 'Nguoi khac han ta khi ho xuc pham ta. Loi cua ta la phan ung cua ta.'\\
\textit{(He wrote to himself: 'Others are beyond my control when they offend me. What is mine is my reaction.')}

Moi sang ong tu nhan minh tru'oc khi ra quyet dinh --- nhu mot triet hoc tan tinh toan mon nhung quan doi kho ung.\\
\textit{(Every morning he cross-examined himself before making decisions --- like a rigorous philosopher commanding difficult armies.)}

Ong lam nhung gi ong viet --- va de che La Ma song ben vung them hai the ky nho tinh than do.\\
\textit{(He did what he wrote --- and the Roman Empire endured two more centuries partly because of that spirit.)}

\end{truyen}
\begin{baihoc}
  Nhat ky tu chon thuc la truong tap luyen su tu chon --- viet ra de kiem tra chinh minh truoc, roi moi hanh dong.\\
  \textit{(A self-examination journal is actually a self-discipline training ground --- write to test yourself first, then act.)}
\end{baihoc}

\section{Churchill --- Loi Noi Cuu Nuoc Anh}
\begin{truyen}{Churchill --- Loi Noi Cuu Nuoc Anh}{Winston Churchill --- Anh 1940}
\chuhoa{T}{}T hang 5/1940, quan Duc da chiem Phap, da pha chong Hoa Lan va Bi --- nuoc Anh gan nhu don doc.\\
\textit{(In May 1940, Germany had taken France, had broken through Holland and Belgium --- Britain stood nearly alone.)}

Churchill len lam Thu tuong trong hoan canh tan nhoai nhat co the tuong tuong; ong len san dien doc bai phat bieu.\\
\textit{(Churchill became Prime Minister in the most desperate circumstances imaginable; he mounted the podium to deliver a speech.)}

Khong mot loi hua chien thang de dang; ong noi: 'Toi khong co gi de dem het ngoai mau, mo hoi va nuoc mat.'\\
\textit{(Not one promise of easy victory; he said: 'I have nothing to offer but blood, toil, tears, and sweat.')}

Nhung Quoc Hoi Anh va ca nuoc Anh cam thay khac --- ho cam thay duoc nghe thay su that lan lanh dao that su.\\
\textit{(But the British Parliament and the whole nation felt something different --- they felt they were hearing the truth and real leadership.)}

Loi noi thieu su that khong bao gio dong vien duoc ai; chi loi noi song voi su that moi co the cam hoa con tim.\\
\textit{(Words devoid of truth never truly inspire anyone; only words alive with truth can move hearts.)}

\end{truyen}
\begin{baihoc}
  Hanh dong trung thuc --- ke ca khi su that la bi thu --- luon hieu qua hon su lua doi duong nhu dep den khi can thiet.\\
  \textit{(Honest action --- even when the truth is bleak --- always works better than comforting deception when it counts.)}
\end{baihoc}

\section{Mark Twain --- Noi Dieu Ban tin, Khong Noi Dieu Ban Muon}
\begin{truyen}{Mark Twain --- Noi Dieu Ban tin, Khong Noi Dieu Ban Muon}{Mark Twain --- Hoa Ky 1880}
\chuhoa{M}{}M ark Twain la nhà van My noi tieng it ai biet rang ong biet cach im lang gioi hon cach noi.\\
\textit{(Mark Twain was a famous American writer who few knew was equally skilled at knowing when to be silent.)}

Mot lan ong duoc moi phat bieu truoc mot dam dong uu tu ve chinh sach thue khoa bat cong.\\
\textit{(Once he was invited to speak before a crowd angry about unjust tax policies.)}

Thay dam dong dang nuong nau, thay vi gap them lua, ong chi noi mot cau: 'Nhung gi mieng tui dang noi het roi.'\\
\textit{(Seeing the crowd already heated, instead of fueling the fire, he said only one sentence: 'The empty pockets say it all.')}

Bai phat bieu ngan gon nhat trong su nghiep ong --- nhung ghi nho lau nhat khi nhung nguoi chung kien ke lai.\\
\textit{(His shortest speech ever --- but remembered the longest by those who witnessed it.)}

Twain viet: 'Su khac biet giua tu dung va tu suyt dung cung lon nhu su khac biet giua sap setam va con doi'.\\
\textit{(Twain wrote: 'The difference between the almost right word and the right word is the difference between the lightning bug and lightning.')}

\end{truyen}
\begin{baihoc}
  Mot tu dung dung co the thay doi cuoc tro chuyen; phan con lai cua ngon ngu la de bao ve tu dung do.\\
  \textit{(One right word can change a conversation; the rest of language exists to protect that right word.)}
\end{baihoc}

\section{Ho Xuan Huong --- Tho La Chien Luoc Xa Hoi}
\begin{truyen}{Ho Xuan Huong --- Tho La Chien Luoc Xa Hoi}{Ho Xuan Huong --- Viet Nam TK 18}
\chuhoa{H}{}H o Xuan Huong song trong xa hoi phong kien khe khe rang buoc phu nu --- nhung ba da su dung tho ca nhu vu khi.\\
\textit{(Ho Xuan Huong lived in a feudal society that tightly restricted women --- but she used poetry as a weapon.)}

Ba khong can an thuong, khong can xin phep --- ba noi dieu ba muon noi qua cac bai tho day an du.\\
\textit{(She did not need permission or approval --- she said what she wanted through poems full of metaphors.)}

Nhung bai tho cua ba co hai tang y nghia: tang ngoai von lat va tang trong can dam dot pha.\\
\textit{(Her poems had two layers of meaning: a surface layer of wit and an inner layer of bold subversion.)}

Nhung ke doc tu tuong chi thay cai ben ngoai; nhung ke hieu thi kinh ngac truoc su can dam cua tu tuong.\\
\textit{(Those who read narrowly only saw the surface; those who understood were stunned by the courage of the thought.)}

Ngon ngu --- dung kheo --- la thu vu khi co the xuyen qua rao can xa hoi ma khong can bat ky bao luc nao.\\
\textit{(Language --- used cleverly --- is a weapon that can penetrate social barriers without needing any violence.)}

\end{truyen}
\begin{baihoc}
  Khi ban khong the noi thang, hay hoc cach noi vong quanh --- va doi khi con xuc tich va sac ben hon.\\
  \textit{(When you cannot speak directly, learn to speak indirectly --- and sometimes it is even more precise and sharp.)}
\end{baihoc}

\section{Bertrand Russell --- Hanh Dong Theo Nhung Gi Ban Tin}
\begin{truyen}{Bertrand Russell --- Hanh Dong Theo Nhung Gi Ban Tin}{Bertrand Russell --- Anh TK 20}
\chuhoa{B}{}B ertrand Russell la mot triet hoc gia noi tieng cua nuoc Anh da nghi den chien tranh va hoa binh trong suot cuoc doi.\\
\textit{(Bertrand Russell was a famous British philosopher who thought about war and peace throughout his entire life.)}

Trong the chien I, khi chan tranh duoc ton vinh, ong con lai la mot trong so it nguoi cong khai phan doi.\\
\textit{(During World War I, when war was being glorified, he was among the very few who publicly opposed it.)}

Ong bi mat viec lam o dai hoc, bi bat giam 6 thang --- nhung khong bao gio im lang hay tu choi.\\
\textit{(He lost his university position, was imprisoned for 6 months --- but never went silent or recanted.)}

Gan cuoi doi, o tuoi 90, ong xuong duong bieu tinh chong vu khi hat nhan giua Chien tranh Lanh.\\
\textit{(Near the end of his life at 90, he marched in the streets protesting nuclear weapons during the Cold War.)}

Ong noi: 'Cam giac toi pham nghiem trong nhat la lam dieu gi do ma ban biet la sai vi nguoi khac mong ban lam.'\\
\textit{(He said: 'The most serious moral crime is doing something you know is wrong because others expect you to.')}

\end{truyen}
\begin{baihoc}
  Noi dieu ban tin --- du the gioi cho ban la sai --- la hanh dong ung xu tri thuc cao nhat con nguoi co the thuc hien.\\
  \textit{(Speaking what you believe --- even when the world tells you you're wrong --- is the highest intellectual act a person can perform.)}
\end{baihoc}
